\documentclass{gak-article}

\usepackage{preamble}
\usepackage{pftools}
\usepackage{lipsum}
\usepackage{minted}
\usepackage{maths}
\usepackage{lambda-macros}
\usepackage{macros}

\addbibresource{refs.bib}

\title{A new paper dawns}
\author{G. A. Kavvos \and Alex Kavvotsky}
\date{\today}

\begin{document}

\maketitle

\begin{abstract}
  \lipsum[1]
\end{abstract}

\tableofcontents


\section{First}
  \label{section:first}

\subsection{First one}
  Here is some mathematics: \( \int \frac{1}{x} = \ln |x| + C \).
  
  Here is a display block of mathematics:
  \[
    \mathbf{Y} \defeq (\lambda x. xx)(\lambda x. xx)
  \]

  \lipsum[1] \parencite{qian2021}

\subsection{Second one}
  \lipsum[2]
  
  \paragraph{Another paragraph}
  Here's a grammar:
  \begin{grammar}
    Terms & M, N & x \mid \Lam{x}{M} \mid \App{M}{N}
  \end{grammar}

  \lipsum[3]
  
  \paragraph{Rules}

  Here are some proof rules:
  \begin{mathpar}
    \inferrule{
      \IsTerm[\Gamma, \DeclVar{x}{A}]{M}{B}
    }{
      \IsTerm{\Lam{x}[A]{M}}{A \to B}
    }
    
    \and

    \inferrule{
      \IsTerm{M}{A \to B} \\
      \IsTerm{N}{A}
    }{
      \IsTerm{\App{M}{N}}{B}
    }
  \end{mathpar}
  There are various ways to label these rules: look up the \texttt{mathpartir} package documentation.

  Or, if you load \texttt{pftools}, you can both label \emph{and} link to rules:
  \begin{mathpar}
    \inferH{Lam}{
      \IsTerm[\Gamma, \DeclVar{x}{A}]{M}{B}
    }{
      \IsTerm{\Lam{x}[A]{M}}{A \to B}
    }
    
    \and

    \inferH{App}{
      \IsTerm{M}{A \to B} \\
      \IsTerm{N}{A}
    }{
      \IsTerm{\App{M}{N}}{B}
    }
  \end{mathpar}
  You can then refer to rule \ruleref{Lam}.
  
  \begin{theorem}
      \label{theorem:induction}
    This can be proven by induction.
  \end{theorem}
  
  We now prove \cref{theorem:induction} by induction.
  
  \begin{proof} \hfill
    \begin{indproof}
      \indcase{Base} Base case proof.

      \indcase{Inductive} Inductive case proof.
    \end{indproof}
  \end{proof}

  
\section{Second}

Recall that we proved \cref{theorem:induction} in \cref{section:first}.

Code is typeset with \texttt{minted}:
\begin{minted}{haskell}
data Term = Var String | Lam String Term | App Term Term

free :: Term -> [String]
free (Var x)   = [x]
free (Lam x m) = filter (/= x) (free m)
free (App m n) = free m ++ free n
\end{minted}

\lipsum[4]

\printbibliography

\end{document}
